Настоящее приложение фиксирует стилистические соглашения, принятые авторами при написании книги, и содержит сводку ключевых обозначений и математических команд LaTeX, используемых в тексте.
Владение LaTeX мы считаем не просто полезным, а необходимым навыком для студента, всерьёз занимающегося математикой и смежными науками.
Он пригодится при конспектировании лекций и семинаров, подготовке домашних и курсовых работ, написании дипломной работы и научных статей, а в дальнейшем~--- в любой профессиональной деятельности, связанной с подготовкой технических текстов.
В рамках нашего курса мы стремимся к тому, чтобы студенты не только освоили содержательную часть математического анализа, но и научились грамотно и единообразно записывать свои рассуждения.
Именно поэтому в книгу включено настоящее приложение: оно служит практической памяткой по математическому набору и демонстрирует один из возможных подходов к организации крупного LaTeX-проекта.

При этом мы не ставим задачу привести общее введение в LaTeX.
Базовые сведения об установке компилятора, структуре документа, подключении пакетов, работе с рисунками и таблицами и~т.п. подробно изложены в многочисленных учебных пособиях и онлайн-ресурсах, и соответствующий материал читатель легко найдёт самостоятельно.
Здесь же мы сосредотачиваемся на том, что специфично именно для математики: выборе обозначений, унифицированном оформлении определений и утверждений, типографских тонкостях набора формул и~т.п.
Читатель, желающий использовать данное оформление как отправную точку для собственных LaTeX-документов, может найти минимальный шаблон с подключёнными стилями в нашем репозитории.

Книга компилируется с использованием LaTeX-движка \verb|lualatex|.
Стилевое оформление документа вынесено в четыре файла, находящихся в папке \verb|style/| и подключаемых в основном файле \verb|book.tex|:
\begin{itemize}
    \item \verb|color.tex|~---
        палитра книги: базовые цвета, а также семантические цвета для блоков 
            определений,
            утверждений,
            лемм,
            теорем,
            следствий,
            примеров,
            замечаний,
            алгоритмов и
            программного кода;
    \item \verb|base.tex|~---
        подключаемые пакеты и базовые стили (геометрия страницы, шрифты, колонтитулы, заголовки глав и разделов и~т.д.);
    \item \verb|helpers.tex|~---
        собственные математические команды и текстовые сокращения, систематически описанные в последующих разделах;
    \item \verb|block.tex|~---
        собственные команды-обёртки для блоков
            определений (\verb|\defi|),
            утверждений (\verb|\stat|),
            лемм (\verb|\lemm|),
            теорем (\verb|\theo|),
            следствий (\verb|\conc|),
            примеров (\verb|\exam|),
            замечаний (\verb|\note|),
            алгоритмов (\verb|\algo|) и
            программного кода (\verb|\code|).
\end{itemize}

Далее в разделе <<Общие правила оформления>> мы фиксируем принятые нами соглашения по оформлению текста (правила использования многоточий, таблиц, списков и~т.д.).
Последующие разделы упорядочены по темам, при этом в каждом разделе приводится таблица, в которую включены как собственные команды, так и многие стандартные LaTeX-команды, которые регулярно используются в тексте книги.